\documentclass[a4paper,11pt]{article}

% Core packages 
\usepackage[utf8]{inputenc}
\usepackage[T1]{fontenc}
\usepackage{lmodern}
\usepackage[margin=2.2cm]{geometry}
\usepackage{microtype}
\usepackage{amsmath,amssymb}
\usepackage{graphicx}
\usepackage{hyperref}
\usepackage{enumitem}
\usepackage{tabularx}
\usepackage{xcolor}

% Compact, tidy formatting 
\usepackage{titlesec}
\titlespacing*{\section}{0pt}{0.6em}{0.3em}
\titlespacing*{\subsection}{0pt}{0.5em}{0.2em}
\setlist{itemsep=0.2em, topsep=0.2em}
\setlength{\parskip}{0.25em}
\setlength{\parindent}{0pt}

% Gantt chart 
\usepackage{pgfgantt}

% Define lengths only if not already defined
\makeatletter
\@ifundefined{glabelwidth}{\newlength{\glabelwidth}}{}
\@ifundefined{gchartwidth}{\newlength{\gchartwidth}}{}
\makeatother

% Fixed label column and macro to recompute chart width per \linewidth
\setlength{\glabelwidth}{4cm}                         % label column width
\newcommand{\resetganttwidth}{%
  \setlength{\gchartwidth}{\dimexpr\linewidth-\glabelwidth\relax}%
  \ganttset{expand chart=\gchartwidth}%
}
\resetganttwidth

\newcount\ganttweekstart
\newcount\ganttweekend
\newcommand{\weekspanlabel}[2]{%
  \pgfcalendardatetojulian{#1}{\ganttweekstart}%
  \pgfcalendardatetojulian{#2}{\ganttweekend}%
  \advance\ganttweekend by -\ganttweekstart
  \advance\ganttweekend by 1\relax
  \gantttitle{Week}{\the\ganttweekend}%
}
\pdfstringdefDisableCommands{%
  \def\weekspanlabel#1#2{Week}%
  \def\({}%
  \def\){}%
  \def\leftrightarrow{<->}%
}

% Colours
\definecolor{S2}{HTML}{6D28D9}   % Semester 2 accent
\definecolor{POST}{HTML}{6D28D9} % Post-S2 tint
\definecolor{MS}{HTML}{D97706}   % Milestone colour

\ganttset{calendar week text=\currentweek}

% Global pgfgantt style  (note % at EOL to avoid linebreak issues)
\ganttset{%
  expand chart=\gchartwidth,%
  y unit title=0.70cm,%
  y unit chart=0.48cm,%
  % Left label column (fixed width -> identical left edge across charts)
  bar label font=\footnotesize,%
  milestone label font=\footnotesize\bfseries,%
  group label font=\footnotesize\bfseries,%
  bar label node/.append style={text width=\glabelwidth,align=right},%
  milestone label node/.append style={text width=\glabelwidth,align=right},%
  group label node/.append style={text width=\glabelwidth,align=right},%
  % Draw only weekly vertical rules (hide 6 day rules, draw the 7th)
  vgrid={*6{draw=none},*1{draw=black!28}},%
  hgrid=false,%
  % Make milestones pop
  milestone/.append style={fill=MS,draw=MS,line width=0.9pt},%
  milestone height=.85,%
  % Neutral bars + slim title banners
  bar/.append style={fill=black!12},%
  group/.append style={fill=S1!35,draw=S1!80},%
  group top shift=0.15,%
  group height=0.50,%
  title/.append style={draw=black!35,fill=black!6,line width=0.3pt},%
  progress label text={},%
}

% Metadata
\title{Project Outline: Hexapod Robot}
\author{Archie James \quad | \quad Supervisor: Michael O'Toole}
\date{\today}

\begin{document}
\maketitle
\vspace{-1.2em}

\section{Background}
Wheeled Robots struggle over complex terrain where Legged Robots thrive. I aim to develop my skills in Legged Robotics by building on my prior work on an Arduino-based hexapod. While a good start it had no battery charging circuitry and it used cheap servo motors with off the shelf microcontroller boards. This project aims to develop a smaller Hexapod Robot with custom designed PCBs, closed-loop position control of joints, a dedicated Battery Management System and an improved walking algorithm that takes terrain into account and has path planning. I've also already made a bespoke simulation that makes use of Inverse Kinematics and a custom walking algorithm and plan to carry that over into this project.

The project will start by performing tests on a single servo motor, modifying it to provide closed-loop position control. Then I will develop a terrain generator and path planner in simulation to improve the robot's ability to navigate rough terrain. Finally, I will design a battery management system to allow for untethered operation and integrate all components into a final demonstration.

\section{Overall Aim}
To deliver a hexapod robotics stack that will bridge simulation and hardware: (i) a closed-loop servo module, (ii) a terrain suite for parameterised rough terrain, and (iii) a safe, rechargeable 2-cell battery subsystem. The project should result in a single-leg rig and eventually a full chassis that demonstrate behaviours developed in simulation. A path planner may also be explored as a stretch objective, but the core focus is on these deliverables and their evidence.


\section{Objectives (Assessable Deliverables)}
\begin{enumerate}[nosep,label=\textbf{O\arabic*}]
  \item Closed-Loop Servo (STM32) - bench-validated tracking on a single-leg rig; steady-state accuracy and settling time within thresholds suitable for a stable gait. Final numeric targets will be set after initial bring-up and tests and recorded in the logbook.
  \item Reproducible Terrain \& Benchmark Suite - parameterised terrains with a system to report success, slip, and stability metrics on fixed scenarios. Any performance thresholds will be fixed after the first simulation experiments.
  \item Path Planner (stretch) - scope and acceptance criteria defined after further research.
  \item 2S Battery System (USB-C PD, Charger, Protection, Balancing) - PD negotiation demonstrated; charger operation validated; primary protection trips; balancing limits inter-cell drift within an agreed tolerance. Exact tolerances will be aligned to datasheets and set following initial charge/discharge tests.
  \item Integration Demo (Sim \(\leftrightarrow\) Hardware) - comparable behaviour between simulator outputs and hardware (single-leg then full chassis) on a defined mini-suite of tasks. Criteria to be set when servo and terrain baselines are established.
\end{enumerate}

% --- Shared widths & colours for both charts ---
% \newlength{\glabelwidth}
\setlength{\glabelwidth}{4cm} % left label column (same for both charts)
\resetganttwidth

\definecolor{S1}{HTML}{0E7490}      % teal-ish
\definecolor{S2}{HTML}{6D28D9}      % purple-ish
\definecolor{Post}{HTML}{6D28D9}    % reuse S2 hue
\definecolor{MS}{HTML}{D97706}      % milestone colour (amber)

% Week labels: "W1 W2 …"



% \section{Provisional Project Plan}

\begin{center}
\begin{minipage}{\linewidth}
\centering

% SEMESTER 1 
\medskip
\makebox[\linewidth]{\Large\bfseries Semester 1 (Weeks 4-12)}
\resetganttwidth
\begin{ganttchart}[
  time slot unit=day,
  time slot format=isodate
]{2025-10-20}{2025-12-21} % Monday W4 -> Sunday end of W12

  % Title row 1: months (aligned to chart span)
  \gantttitlecalendar*[
    title list options={
      title/.append style={draw=black!45,line width=0.35pt,fill=black!6},
      title label node/.append style={inner sep=1.2pt}
    }
  ]{2025-10-20}{2025-12-21}{month=shortname} \\
  \weekspanlabel{2025-10-20}{2025-12-21} \\
  % Title row 2: week numbers (W4-W12 span)
  \gantttitlecalendar*[
    title list options={
      title/.append style={draw=black!45,line width=0.35pt,fill=black!6},
      title label node/.append style={inner sep=1.2pt}
    }
  ]{2025-10-20}{2025-12-21}{week=4} \\

  % Milestones tied to Semester 1
  \ganttmilestone{PRM1}{2025-10-27} \\
  \ganttmilestone{PRM2}{2025-11-24} \\
  \ganttmilestone{Draft Intro (2pm)}{2025-12-12} \\

  % Servo-focused kick-off (Weeks 4-6)
  \ganttnewline
  \ganttgroup{Servo}{2025-10-20}{2025-11-09} \\
  \ganttbar{Hardware mods \& sensing}{2025-10-20}{2025-10-31} \\
  \ganttbar{Control loop tuning}{2025-10-29}{2025-11-09} \\
  \ganttmilestone{Servo demo ready}{2025-11-08} \\

  % Simulation brought in line (Weeks 7-9)
  \ganttnewline
  \ganttgroup{Simulation Alignment}{2025-11-10}{2025-11-30} \\
  \ganttbar{Refactor}{2025-11-10}{2025-11-20} \\
  \ganttbar{Hardware \(\leftrightarrow\) sim bridge}{2025-11-18}{2025-11-30} \\
  \ganttmilestone{Sim bridge OK}{2025-11-29} \\

  % Terrain generator push (Weeks 10-12)
  \ganttnewline
  \ganttgroup{Terrain Generator}{2025-12-01}{2025-12-21} \\
  \ganttbar{Heightmap pipeline}{2025-12-01}{2025-12-12} \\
  \ganttbar{Metrics \& validation}{2025-12-08}{2025-12-21} \\
  \ganttmilestone{Terrain demo?}{2025-12-20}
\end{ganttchart}

% SEMESTER 2
\medskip
\makebox[\linewidth]{\Large\bfseries Semester 2 (Weeks 1-12)}
\resetganttwidth
\begin{ganttchart}[
  time slot unit=day,
  time slot format=isodate,
  title/.append style={fill=S2!10}, % S2 header tint
  group/.append style={fill=S2!30,draw=S2!80}
]{2026-02-02}{2026-05-17} % W1 Monday -> Sunday after W15

  \gantttitlecalendar[
    title list options={
      title/.append style={draw=black!45,line width=0.35pt,fill=black!6},
      title label node/.append style={inner sep=1.2pt}
    }
  ]{month=shortname} \\
  % Show W1..W12
  \weekspanlabel{2026-02-02}{2026-04-26} \\
  \gantttitlecalendar*[
    title list options={
      title/.append style={draw=black!45,line width=0.35pt,fill=black!6},
      title label node/.append style={inner sep=1.2pt}
    }
  ]{2026-02-02}{2026-04-26}{week=1} \\

  % Reviews & deliverables
  \ganttmilestone{PRM3}{2026-02-16} \\
  \ganttmilestone{Draft Final Report}{2026-03-13} \\
  \ganttmilestone{PRM4}{2026-03-16} \\
  \ganttmilestone{Final Report}{2026-05-01} \\

  % Battery management boot-up (Weeks 1-3)
  \ganttnewline
  \ganttgroup{Battery Mgmt}{2026-02-02}{2026-02-23} \\
  \ganttbar{PD negotiation \& charger}{2026-02-02}{2026-02-16} \\
  \ganttbar{Protection \& balancing tests}{2026-02-09}{2026-02-23} \\
  \ganttmilestone{Battery Ready}{2026-02-23} \\

  % Integration thrust (Weeks 4-15)
  \ganttnewline
  \ganttgroup{Integration}{2026-02-23}{2026-05-17} \\
  \ganttbar{Single leg validation}{2026-02-23}{2026-03-15} \\
  \ganttbar{Full chassis + BMS}{2026-03-16}{2026-04-12} \\
  \ganttbar{System tuning \& trials}{2026-04-06}{2026-05-04} \\
  \ganttbar{Demo prep \& window}{2026-05-05}{2026-05-17} \\
  \ganttmilestone{Final Demo}{2026-05-15}
\end{ganttchart}

\end{minipage}
\end{center}

\end{document}
